\section{Variables y operacionalización}\label{variables-operacionalizacion}

\subsection{Variables de investigación}\label{variables-de-investigaciuxf3n}

\subsubsection{Variable independiente}\label{variable-independiente}

\textbf{Condición de adaptación del modelo.}

Se trata de una variable categórica con dos niveles:

\begin{itemize}
\tightlist
\item
  A: modelo original.
\item
  B: mismo modelo adaptado mediante QLoRA.
\end{itemize}

\subsubsection{Variable dependiente}\label{variable-dependiente}

\textbf{Calidad de las respuestas breves a preguntas cotidianas en quechua sureño.}

Se medirá mediante pertinencia, corrección lingüística y adecuación a la variedad.

El indicador principal será la aceptabilidad global de la respuesta.

\subsubsection{Variables de control y estratificación}\label{variables-de-control-y-estratificaciuxf3n}

Se controlarán:

\begin{itemize}
\tightlist
\item
  Identidad y versión del modelo base.
\item
  Conjunto de preguntas de prueba.
\item
  Plantilla de instrucciones.
\item
  Configuración de generación.
\item
  Límite máximo de tokens de salida.
\item
  Entorno de inferencia.
\item
  Criterios de evaluación.
\end{itemize}

Se registrarán como factores de estratificación:

\begin{itemize}
\tightlist
\item
  Variedad lingüística.
\item
  Dominio temático.
\item
  Tipo de pregunta.
\item
  Longitud de entrada.
\end{itemize}

\subsection{Operacionalización de variables}\label{operacionalizaciuxf3n-de-variables}

La Tabla~\ref{tab:operacionalizacion} relaciona cada variable con sus dimensiones, indicadores, escala de medición y fuente prevista.

\begin{table}[H]
\centering
\begin{threeparttable}
\caption{Operacionalización de las variables}
\label{tab:operacionalizacion}
\small
\begin{tabularx}{\textwidth}{@{}P{2.45cm}P{2.4cm}L P{1.7cm}P{2.4cm}@{}}
\toprule
\textbf{Variable} & \textbf{Dimensión} & \textbf{Indicador} & \textbf{Escala} & \textbf{Fuente} \\
\midrule
Condición de adaptación & Estado del modelo & Modelo original o el mismo modelo adaptado mediante QLoRA & Nominal & Configuración experimental \\
Calidad de respuesta & Pertinencia & Puntuación de la rúbrica: 0, 1 o 2 & Ordinal & Evaluación humana \\
Calidad de respuesta & Corrección lingüística & Puntuación de la rúbrica: 0, 1 o 2 & Ordinal & Evaluación humana \\
Calidad de respuesta & Adecuación a la variedad & Puntuación de la rúbrica: 0, 1 o 2 & Ordinal & Evaluación humana \\
Calidad de respuesta & Aceptabilidad global & Aceptable o no aceptable según regla predefinida & Nominal dicotómica & Rúbrica adjudicada \\
Desempeño técnico complementario & Memoria & Pico de memoria de GPU durante entrenamiento e inferencia & Razón & Registro técnico \\
Desempeño técnico complementario & Latencia & Segundos por respuesta & Razón & Registro técnico \\
Desempeño técnico complementario & Entrenamiento & Duración y número de parámetros entrenables & Razón & Registro técnico \\
\bottomrule
\end{tabularx}
\begin{tablenotes}
\item \textit{Nota.} Las medidas técnicas serán descriptivas y no constituirán una segunda variable dependiente principal.
\end{tablenotes}
\end{threeparttable}
\end{table}

\subsection{Criterio de aceptabilidad global}\label{criterio-de-aceptabilidad-global}

Cada dimensión recibirá una puntuación de 0 a 2.

Una respuesta será clasificada como aceptable si cumple simultáneamente:

\begin{itemize}
\tightlist
\item
  Pertinencia: 2 puntos.
\item
  Corrección lingüística: al menos 1 punto.
\item
  Adecuación a la variedad: al menos 1 punto.
\end{itemize}

Una respuesta que no conteste la pregunta no será aceptable, aunque su formulación sea lingüísticamente correcta.

Este criterio deberá ser aprobado durante la validación de la rúbrica, antes de la evaluación final.

Para cada condición $c$ y variedad, la proporción de respuestas aceptables se calculará como:
\begin{equation}
\widehat{p}_{c}=\frac{n_{\mathrm{aceptables},c}}{N_c},
\label{eq:proporcion_aceptabilidad}
\end{equation}
donde $n_{\mathrm{aceptables},c}$ es el número de respuestas clasificadas como aceptables y $N_c$ es el número de preguntas de prueba de esa variedad evaluadas en la condición $c$.

La diferencia pareada de interés entre las dos condiciones se expresará como:
\begin{equation}
\widehat{\Delta p}=\widehat{p}_{\mathrm{QLoRA}}-\widehat{p}_{\mathrm{original}}.
\label{eq:efecto_adaptacion}
\end{equation}
El signo y la magnitud se estimarán con los datos; no se presupone una mejora.

\subsection{Matriz de consistencia}\label{matriz-de-consistencia}

La Tabla~\ref{tab:matriz-consistencia-capitulo} muestra la correspondencia entre los problemas, los objetivos, las hipótesis y el método propuesto.

\begin{table}[H]
\centering
\begin{threeparttable}
\caption{Matriz de consistencia de la investigación}
\label{tab:matriz-consistencia-capitulo}
\scriptsize
\begin{tabularx}{\textwidth}{@{}P{3.45cm}P{3.7cm}P{2.65cm}L@{}}
\toprule
\textbf{Problema} & \textbf{Objetivo} & \textbf{Variable o resultado} & \textbf{Procedimiento} \\
\midrule
¿Qué efecto tiene adaptar mediante QLoRA un modelo de lenguaje compacto sobre la calidad de sus respuestas breves a preguntas cotidianas en quechua sureño? & Evaluar el efecto de adaptar mediante QLoRA un modelo de lenguaje compacto sobre la calidad de sus respuestas breves a preguntas cotidianas en quechua sureño. & Condición de adaptación; calidad de respuestas y aceptabilidad global. & Comparación pareada con las mismas preguntas; evaluación humana ciega y análisis por variedad. \\
\addlinespace[0.15cm]
¿Cómo construir y validar un conjunto de preguntas y respuestas cotidianas que permita evaluar las variedades chanka y collao sin confundir datos de traducción con datos de generación de respuestas? & Construir y validar un corpus de preguntas y respuestas breves de uso cotidiano en quechua sureño, con identificación de las variedades chanka y collao y separación de datos para entrenamiento, validación y prueba. & Calidad y composición del corpus & Auditoría de fuentes, revisión lingüística y partición por grupos semánticos. \\
\addlinespace[0.15cm]
¿Cómo adaptar un único modelo de lenguaje compacto mediante QLoRA para la generación de respuestas breves en quechua sureño, considerando las restricciones computacionales disponibles? & Evaluar la factibilidad computacional y, si resulta viable, adaptar un único modelo de lenguaje compacto mediante QLoRA y documentar su configuración. & Condición de adaptación & Piloto, ajuste supervisado y registro de configuración si el entrenamiento es viable. \\
\addlinespace[0.15cm]
¿Qué diferencias se observan en la aceptabilidad, pertinencia, corrección lingüística y adecuación a la variedad de las respuestas del modelo antes y después de la adaptación? & Comparar, mediante evaluación humana y análisis estadístico pareado, la calidad de las respuestas generadas por el mismo modelo antes y después de su adaptación, presentando resultados generales y diferenciados por variedad. & Aceptabilidad y dimensiones de calidad & Generación controlada, evaluación ciega con rúbrica y análisis estadístico pareado. \\
\bottomrule
\end{tabularx}
\end{threeparttable}
\end{table}

\begin{center}\rule{0.5\linewidth}{0.5pt}\end{center}
